\section{Verification of the numeral proof constructors}
\label{sec:annex-p05-pa-numeral-proof-builders}

\begin{annexitem}{annex:p5.builder-nocapture}{thm:p05-builder-nocapture}{The specified substitution machine on the required inputs}
\begin{proof}
All arguments in this annex take place in \(\PA\). Each named function
means the unique graph output of its fixed description. The graph equations,
table-prefix equations, constructor decoding and list equations used below
are those proved in \cref{annex:p5.checker-description}; induction is full
\(\PA\) induction with the other displayed parameters retained. In
particular, the induction arguments apply to the numerical variables of
\(\PA\), rather than just to standard codes.

The new table is a finite description. Its row update uses only the previous
table, the fixed type/variable tests, projections, constructors, arithmetic
and a finite case distinction. Its bound is \(Sc\), evaluated before the
fold. Every valid expression child has smaller code, so its row is available.
Invalid input has exactly the displayed three zero fields. The earlier
table-prefix theorem identifies a retrieved child row with its own completed
table value, internally and on all codes.

Strong induction on \(c\), expressed by retaining all earlier instances,
proves the following assertion whenever its first field is one: the returned
expression has the input type, and the cost is positive and at most twice
the input node count. At a variable the replacement has the required term
type, or the variable is unchanged. Zero and bottom retain their types.
At a unary or binary constructor the child assertions and the fixed grammar
give the output type. The cost is two plus the child costs, at most twice
one plus the sum of child node counts. At a stopped binder the unchanged
expression has cost one. At another binder the matrix assertion and the
same cost calculation apply. These exhaust the valid syntax cases.

Here is a stronger execution invariant identifying the actual selected
machine, including its time bound. Write \(\mathsf{step}\) for its fixed
single update and \(\mathsf{step}^{,w}\) for its \(w\)-fold primitive
recursion. For arbitrary raw continuation and value lists \(D,V\), if
the table row at \(c\) is \([1,z,w]\), then
\[
 \mathsf{step}^{,w}([1,C([0,c,i,u],D),V])=[1,D,C(z,V)].
\tag{1}
\]
This is an arithmetic equation of specified graph outputs. Prove it in the
same strong induction on \(c\), with \(D,V\) universally quantified.
At a variable, zero or bottom the first substitution command returns the
declared value in one update. At a nonbinding unary node, one update puts
its child command and assembly command before \(D\). The child induction
instance returns its value, and the one assembly update wraps the specified
tag. At a binary node, the first update puts left command, right command,
and assembly in that order. Apply the left instance with the other commands
as continuation, then the right instance with the returned left value on
the stack. Assembly pops right and then left, giving exactly
\(N_k(z_0,z_1)\) in \(2+w_0+w_1\) updates. The iteration equation
\(\mathsf{step}^{a+b}(s)=\mathsf{step}^{b}(\mathsf{step}^{a}(s))\),
proved by induction on \(b\), justifies these compositions of runs.

At a binder with index \(i\), the actual machine stops and returns \(c\)
in one update. At a different binder \(j\), table success gives
\(\mathsf{AV}(u,j)=0\). Thus the actual machine selects its no-renaming
branch, invokes the matrix command, and then its quantified assembly command.
The matrix induction instance gives (1) in \(2+w_0\) updates. Its term and
formula tests succeed by the previously established type assertion. These
cases verify every machine update that can occur under the table-success
premise. In particular no fresh-renaming or continuation-after-renaming
command is used. The new child values provide every value-stack entry needed
by assembly. The continuation \(D\) is never executed in (1).

Put \(D=V=0\). At cost \(w\) the state is successful with no pending
command and exactly one value. This is an absorbing completed state of the
selected machine. The internally proved bound
\(w\leq2\mathsf{nodes}(c)<20(\mathsf{nodes}(c)+1)^2\) therefore makes
its final extraction return exactly \(z\). This proves
\(\mathsf{sub}=\mathsf{ns}\) from the defining equations of that machine.

If the replacement \(u\) is a term with no variable occurrences, every factor
\(Z(\mathsf{AV}(u,j))\) is one; the same strong induction proves success
at every valid input. For a quantifier-free input there is no such factor,
so every term replacement succeeds. These arguments also give the explicit
nonbinding substitution equations, not just success of a computation.

Induction on the recurrence for \(\nu\), using the syntax-table equations,
proves that it is a term with every all-variable and free-variable bit zero.
Retain \(\mathsf{Tm}(u)=1\) as a parameter while inducting on that
recurrence in the new table. At zero this hypothesis passes the replacement
type test, so the zero-term case gives success bit one and result \(\nu(0)\).
At the successor the same type test passes, and the predecessor induction
instance gives success bit one and the unchanged successor result. Thus
\(\mathsf{nc}(\nu(n),i,u)=1\) and
\(\mathsf{ns}(\nu(n),i,u)=\nu(n)\) under the retained hypothesis.
Equation (1) then gives \(\mathsf{sub}(\nu(n),i,u)=\nu(n)\) under that
same hypothesis, proving the body's universally closed guarded assertion.

For an invalid replacement, \(\mathsf{Tm}(u)=0\), the original type guards
instead give \(\mathsf{nc}(c,i,u)=\mathsf{ns}(c,i,u)=
\mathsf{cost}(c,i,u)=0\) and \(\mathsf{sub}(c,i,u)=0\) for every \(c\).
For example, the numerical input \(u=0\) is an empty-list code, not the
zero-term code \(\nu(0)=N_1()=C(1,0)=2\). At \(n=i=u=0\), both
substitution outputs are therefore zero, while \(\nu(0)=2\). These
unchanged failure values explain the replacement-term hypothesis.

Induction on \(j\) gives
\(\mathsf{it}(s(t),j)=s(\mathsf{it}(t,j))=\mathsf{it}(t,Sj)\):
the base is the zero equation and the step applies one further successor
constructor to the preceding equality.

One freshness fact will be used repeatedly. Strong induction on the syntax
code, with the maximum-index table equations, proves internally that a valid
expression has \(\mathsf{AV}(c,j)=0\) whenever \(j>m(c)\). At a variable
the maximum is its index; at a compound the maximum dominates each child
maximum; at a binder it also dominates the binder. Zero and bottom have no
variable bits. Thus every hole index chosen in the body is absent from the
other terms and binders in its context. The displayed open-term instances
of the fixed axioms have no binding conflict: first substituting for index
0 leaves at most the index-1 binder, and the replacement then has no index 1;
the second substitution is in a quantifier-free matrix. Their table equations
therefore identify all their literal roots. This completes every claimed
substitution prerequisite.
\end{proof}
\end{annexitem}

\begin{annexitem}{annex:p5.builder-row-templates}{thm:p05-builder-row-templates}{Internal tree, row and occurrence equations for the templates}
\begin{proof}
Work with the exact \(\mathsf{Rows}\), \(\mathsf{Tree}\), parent-count
and open-list descriptions. For clarity, a successful tree assertion means
\(\mathsf{Tree}(R)=1\); this includes nonemptiness, every row bit equal
to one, and the exact parent-count product. It does not yet impose a
sentence root or an axiom test on open assumptions.

The joins are fixed finite compositions of length, shift, concatenation and
constructor functions. At arity zero the result is the one displayed row;
at any other allowed arity, positivity of premise lengths puts their root
references below the appended row. The cumulative block intervals are
disjoint. The internal shift and concatenation equations from
\cref{annex:p5.checker-internal-mp} give, for block \(j\) and \(v<l_j\),
\[
 f_{\rm joined}(t_j+v)=f_{R_j}(v),\qquad
 \lambda_{\rm joined}(t_j+v)=\lambda_{R_j}(v).
\tag{2}
\]
The parameter fields are unchanged and premise references receive only the
offset \(t_j\).

Induct through each block's row counter. The first block's row computations
are unchanged. For a later block, the induction assertion is that each row
bit is its old bit and its open list is the old list shifted by \(t_j\).
At assumptions the list is \([t_j+v]\); identity introduction has the empty
list. Nondischarge rules use the old formula and parameter data, and
concatenation commutes with shift. All substitution calls have exactly their
old numerical arguments, so their unique outputs are unchanged. Discharge
filters have equal label/formula tests at corresponding entries by (2);
the internal filter/shift equation preserves both their result and their
strict length-decrease test. Universal introduction has the same free-variable
calls on corresponding leaves. Existential elimination has the same matrix,
eigenvariable, major formula and conclusion; after the same minor discharge,
its all-variable calls and its minor-before-major open-list order are
unchanged. Alpha and equality rules likewise have identical syntactic
arguments. These are the old-row cases of the accepted internal joining
proof, whose premises were successful rows, without a sentence-root premise.
Appending a last row changes none of the prefix lookups used in this induction.

Split the defining parent-count sum into the fixed block intervals and the
last row. Cancellation of offsets gives, for every \(v<l_j\),
\[
 \mathsf{pc}(\mathsf J,t_j+v)
       =\mathsf{pc}(R_j,v)+E(v,D(l_j)).
\tag{3}
\]
No other block contributes, since its references are confined to its own
interval. Each last-row premise names precisely its block root, and these
references are distinct. The new last row has count zero. Substituting the
old count zero at each root and one at every nonroot proves the required
new counts. These equations follow by bounded-sum induction and apply to
arities one, two and three. For arity zero there is just one row and no
reference, so the count is zero directly.

The last row's successful local test now completes the tree product. Its
open list is its exact rule update on the shifted premise lists. In
particular, labels that coincide in different copied trees do not change
an old discharge: the old row still consults only its own premise branch.
This proves the general joining assertion, including temporary assumptions
and intermediate roots with free variables.

For the special equality templates use the internal substitution identities
just proved. With the prescribed fresh hole, symmetry's two premises are
\(t=u\) and \(t=t\), and the matrix is \(x_j=t\); tag 19 returns
\(u=t\). Transitivity uses premises \(u=w\) first and \(t=u\) second,
with matrix \(t=x_j\); it returns \(t=w\). Congruence uses \(t=u\)
and the identity \(g(t)=g(t)\), with matrix \(g(t)=g(x_j)\), and returns
\(g(t)=g(u)\). Every identity premise is a tag-18 row with parameter the
actual term \(t\) or \(g(t)\). The term, formula and freshness facts make
all common and additional row tests one. These nondischarge rules append
only empty-context identity subtrees, so they introduce no new open leaves.

The remaining templates explicitly encode tags 1, 7, 8, 9, 10, 13, 14, 15
and 16 with the exact parameter arities in the body. For example, tag 8
consumes the antecedent tree before the implication tree; tag 10 consumes
the positive formula before its negation. Tags 7 and 9 use the supplied
label, filter its specified formula and require a surviving occurrence.
Tag 16 removes the specified minor instance and tests its eigenvariable
against the all-variable sets of the major, conclusion and remaining open
leaves. These conditions are hypotheses of the template invariant, not
automatic consequences of being a numerical tree.

Finally, if the two roots are sentences and their leaves pass the selected
axiom test, the success guard of \(\mathsf{mp}_{19}\) holds. Its displayed
row expression is precisely the arity-two tag-8 join here after packing,
including the antecedent-first order. Thus the accepted internal MP theorem
applies at such calls. The raw join assertion supplies the explicitly proved
extension needed at the intermediate open-formula calls. A sentence root
and an axiom test on every final open leaf always suffice, by direct expansion
of \(\chi_{19}\), to turn one of the successful raw trees into a selected
proof after packing. The empty final open list contributes the product one.
\end{proof}
\end{annexitem}

\begin{annexitem}{annex:p5.builder-left-addition}{thm:p05-builder-left-addition}{Actual induction records for left addition}
\begin{proof}
The one-row \(\mathsf L(q_j,0)\) is an unmarked assumption. Its formula
is one of the seven literal sentence constants, so its leaf test succeeds.
Each subsequent \(\mathsf{FE}\) is one tag-13 row, with the actual binder,
matrix and replacement term in its three parameter fields. The no-conflict
table identifies the instances used here with the following decoded formulas.
This verifies the whole instance chain rather than treating its final formula
as another axiom.

The induction matrix \(m_Z\) has exactly free index 0. Consequently its
canonical induction axiom \(I_Z\) has no parameter prefix and is literally
\[
 \bigl(0+0=0\land
    \forall x_0((0+x_0=x_0)\to(0+Sx_0=Sx_0))\bigr)
                      \to\forall x_0(0+x_0=x_0).
\tag{4}
\]
The computed zero and successor substitutions in (4) follow from the new
table equations. The exact \(\mathsf{IndAx}\) reconstruction therefore
accepts this leaf. Its code is the fixed expression \(I_Z\), not an
unspecified induction instance.

The root of \(Z_0\) is \(0+0=0\), by the instance of \(q_4\) at
zero. In \(Z_s\), the marked assumption is \(0+x_0=x_0\), marker 1,
label 0. Its successor congruence has root \(S(0+x_0)=Sx_0\). The
instance of \(q_5\) at \((0,x_0)\) has root
\(0+Sx_0=S(0+x_0)\). Transitivity therefore gives exactly the step
consequent \(0+Sx_0=Sx_0\). The marked assumption occurs once in its
open list; the other leaves are unmarked \(q_5\) and identity contributes
none. Tag 7 discharges that occurrence nonvacuously and gives the step
implication. Tag 15 generalizes index 0: every remaining open leaf is a
closed arithmetic axiom. Tag 1 combines the base and universally quantified
step. The raw tag-8 join with (4) concludes \(\mathsf{All}(0,m_Z)\).
Every nonlast row is consumed once by the joining lemma. This is exactly
the defined record \(D_Z\).

For \(m_L\), the free indices are exactly 0 and 1. The induction index
is 1, so \(I_L\) is its induction core preceded by precisely the binder
\(\forall x_0\). The base and successor matrices are respectively
\[
 Sx_0+0=S(x_0+0),\qquad
 Sx_0+Sx_1=S(x_0+Sx_1).
\tag{5}
\]
The no-conflict substitution equations yield these literal codes and the
ordered free-index fold yields prefix \([0]\). This checks its canonical
axiom shape and its final direct-axiom factor.

The record \(L_0\) starts with \(Sx_0+0=Sx_0\), obtained from
\(q_4\) at \(Sx_0\). A separate copy of \(q_4\) at \(x_0\),
successor congruence and symmetry gives \(Sx_0=S(x_0+0)\).
Transitivity gives exactly the first formula in (5). The step record
\(L_s\) chains the following three equalities, in that order:
\[
\begin{aligned}
 Sx_0+Sx_1&=S(Sx_0+x_1),\\
 S(Sx_0+x_1)&=S(S(x_0+x_1)),\\
 S(S(x_0+x_1))&=S(x_0+Sx_1).
\end{aligned}
\tag{6}
\]
The first is \(q_5\) at \((Sx_0,x_1)\). The second is the successor
congruence of the one marked \(m_L\) assumption. The third is the
symmetry of successor congruence of \(q_5\) at \((x_0,x_1)\).
Thus (6) has exactly the step-consequent code in (5). The marker-1
assumption survives once; tag 7 discharges label 0 nonvacuously, and tag 15
generalizes index 1 from closed unmarked axiom leaves. Combine this with
\(L_0\) by tag 1. The other premise is the actual \(\forall E\)
instance of \(I_L\) at term \(v_0\), whose no-conflict identity substitution
returns its induction core literally. Raw implication elimination gives
\(\forall x_1 m_L\); universal introduction at index 0 gives the displayed
root of \(D_L\).

Every operation in (4)--(6) is the fixed template expression in the body,
so these calculations identify actual records, including all references,
parameters and labels after the defined shifts. They are a finite proof
of their checker assertions using its equations and the joining invariant.
The only surviving leaves are among the seven fixed constants and
\(I_Z,I_L\), all unmarked. Direct decoding shows that these nine sentences
have no variable index other than 0 or 1. The syntax and axiom descriptions
on this finite list verify the assertion in \(\PA\) by finitely many
defining-equation instances. No derived conclusion of \(D_Z\) or \(D_L\)
is used as an axiom leaf. Copying either record for a later instance retains
its original verified leaves and adds only its specified \(\forall E\) rows.
\end{proof}
\end{annexitem}

\begin{annexitem}{annex:p5.numeral-arithmetic-builders}{thm:p05-numeral-arithmetic-builders}{Induction on the arithmetic proof recurrences}
\begin{proof}
Use the following arithmetic invariant for the raw trees:
\[
\begin{split}
 \mathsf{Safe}(R,c)\quad\text{means}\quad&
 \mathsf{Tree}(R)=1\land f(R)=c\\
 &\land\forall j<\ell(\mathsf{open}(R))\,
   \bigl(\lambda_R(\mathsf{entry}(\mathsf{open}(R),j))=0
    \land\mathsf{mem}(f_R(\mathsf{entry}(\mathsf{open}(R),j)),
                                             \mathcal A_*)=1\bigr).
\end{split}
\tag{7}
\]
Every operation in (7) is an already specified graph, and the fixed list
\(\mathcal A_*\) was verified in the preceding item. Thus (7) is an
arithmetic formula licensed for induction. If its root is a sentence,
expansion of the selected checker proves
\(\mathsf{Prf}^{19}_{\PA}(\mathsf{pack}(R),c)\). No bound on the
numerical magnitude of \(R\) is needed for this implication.

For \(R_+\), induct on \(n\) in (7) with root
\(\mathsf{Eq}(a(\nu(m),\nu(n)),\nu(A(m,n)))\), retaining \(m\) as parameter.
The base \(q_4\) instance gives exactly
\(\bar m+0=\bar m\), and the equations \(A(m,0)=m\) and \(\nu(0)=o\)
identify the target code. In the step the \(q_5\) instance has root
\[
 \mathsf{Eq}(a(\nu(m),s(\nu(n))),s(a(\nu(m),\nu(n)))).
\]
The preceding proof's successor congruence has root
\(\mathsf{Eq}(s(a(\nu(m),\nu(n))),s(\nu(A(m,n))))\).
Transitivity joins these at identical intermediate term codes. The equations
\(\nu(Sn)=s(\nu(n))\), \(A(m,Sn)=S(A(m,n))\) and
\(\nu(SA(m,n))=s(\nu(A(m,n)))\) identify the exact required new root.
All retained leaves satisfy (7), by the row-joining invariant and the fixed
\(q_5\) leaf check. This is an internal base and step proof, so
\cref{cor:p05-pa-recursion-invariant} proves (7) for the specified recurrence.

For \(R_\times\), use the same induction with root
\(\mathsf{Eq}(b(\nu(m),\nu(n)),\nu(M(m,n)))\). At zero the \(q_6\) instance
and \(M(m,0)=0\) prove it. At the successor, the \(q_7\) instance is
\[
 \bar m\times S\bar n=(\bar m\times\bar n)+\bar m.
\]
The old multiplication proof lifted through right addition by \(\bar m\)
gives
\((\bar m\times\bar n)+\bar m=\overline{M(m,n)}+\bar m\).
The already verified addition builder, at numerical arguments \(M(m,n),m\),
gives this last term equal to \(\overline{A(M(m,n),m)}\). Two transitivity
joins give the stated record's root. The numerical recurrence
\(M(m,Sn)=A(M(m,n),m)\) identifies its last numeral literally. Each
premise tree is copied into its own interval, so no equality proof is shared
as a row with two parents. The leaf invariant persists. These base and step
proofs establish its universal \(\mathsf{Safe}\) assertion in \(\PA\).

For \(R_{\rm right}(t,n)\), retain the permitted term-code parameter
\(t\). The base is the \(q_4\) instance \(t+0=t\), and the step
chains the \(q_5\) instance \(t+S\bar n=S(t+\bar n)\) with successor
congruence of its predecessor. The iterate equation identifies the output
as \(\mathsf{it}(t,Sn)\). For \(R_{\rm left}(t,n)\), the base is
the actual instance of \(D_Z\), \(0+t=t\). The step uses the actual
instance of \(D_L\), \(S\bar n+t=S(\bar n+t)\), and successor
congruence of its predecessor. All instance substitutions satisfy the new
table test: a numeral has no variables, and the only possible open index
of \(t\) is at least 3, away from the index-1 inner binder. Consequently
the displayed formulas are the exact root codes, not merely equivalent
formulas. The same induction proves (7) in both orientations.

The roots of the numeral addition and multiplication builders are sentences,
because numerals have no variables and their constructors add no binders or
variable occurrences. Applying the implication following (7) proves the
two selected-checker assertions in the body. At every use, replacing the
function names by existentially quantified graph outputs gives an ordinary
\(\PA\) derivation. Totality supplies those outputs, their defining
equations supply the base and step identities, and uniqueness identifies
the composed values. This establishes the assertions for the specified
graphs rather than for an extensionally similar representative.
\end{proof}
\end{annexitem}

\begin{annexitem}{annex:p5.numeral-literal-builders}{thm:p05-numeral-literal-builders}{Successor peeling, witness exclusion and literal roots}
\begin{proof}
First verify the raw contradiction constructor in \(\PA\). Assume
\(\mathsf{Tree}(R)=1\), \(k>b\), and
\(f(R)=\mathsf{Eq}(\mathsf{it}(t,k),\nu(b))\), with the stated domain
of \(t\).
Induct on \(j\leq b\) with the invariant that \(P(t,k,b,R,j)\) is a
successful tree with root
\[
 \mathsf{Eq}(\mathsf{it}(t,L(k,j)),\nu(L(b,j))),
\tag{8}
\]
and with exactly the old open leaves of \(R\), shifted as its copied tree,
together with unmarked \(q_2\) leaves. This open-list assertion can be
written using the join's exact shift/concatenation equations; in particular
every old marked leaf survives with its multiplicity.

At zero, subtraction's zero equation gives (8) and the record is \(R\).
At \(j<b\), the arithmetic order and subtraction equations give
\(L(k,j)>0\), \(L(b,j)>0\), and
\[
 L(k,j)=S(L(k,Sj)),\qquad L(b,j)=S(L(b,Sj)).
\]
The \(q_2\) instance in the update has antecedent exactly (8), by the
iterate and numeral equations, and consequent exactly its instance at
\(Sj\). The raw tag-8 join verifies the new tree and open list. Its
term substitutions are numeral terms or successor iterates of \(v_y\)
with \(y\geq3\), so the surviving binder-1 test succeeds and no renaming
changes these literal roots. This proves the step. At \(j\geq b\) the
description retains its state; this defines its total value without adding
any unbounded loop. The final value is evaluated at the fixed bound \(b\).

At \(j=b\), (8) is
\(\mathsf{Eq}(\mathsf{it}(t,L(k,b)),o)\), and \(L(k,b)>0\). Thus the first
term is literally \(s(\mathsf{it}(t,D(L(k,b))))\). The specified
\(q_1\) instance is the negation of precisely this equality. Tag 10,
positive equality first, produces bottom and retains the asserted open
leaves together with one unmarked \(q_1\) leaf. This proves the full
\(\mathsf{Clash}\) invariant. At \(b=0\) the peeling recursion performs
no update; the initial equality is contradicted directly by \(q_1\).

For \(R_{\ne}(m,n)\), internal trichotomy under \(m\ne n\) selects
exactly one of its first two cases. In the first case \(t=o,k=m,b=n\),
and the iterate equals the corresponding numeral. In the second case the
explicit symmetry template first changes the root to
\(\mathsf{Eq}(\nu(n),\nu(m))\); then \(k=n,b=m\). The contradiction invariant
applies in either case. Its only marked leaf is the original equality
\(c\) with marker 1, and it survives once: every update uses its old tree
once, symmetry's identity has no assumptions, and new axiom leaves are
unmarked. Tag 9 with parameter \([0]\) therefore discharges that leaf
nonvacuously and returns exactly \(\mathsf{Neg}(c)\). The resulting tree satisfies
\(\mathsf{Safe}\). If \(m=n\), the equality builder's conditional returns
the premise-free identity tree; equality of the numerical arguments gives
equality of their unique numeral-code outputs, so its root is the displayed
\(\mathsf{Eq}(\nu(m),\nu(n))\). All these roots are sentences.

For the positive weak comparison, \(m\leq n\) gives
\(A(L(n,m),m)=n\). The addition proof therefore has root
\(\mathsf{Eq}(a(\nu(L(n,m)),\nu(m)),\nu(n))\). This is the exact result of
substituting the witness \(\nu(L(n,m))\) into the matrix
\(\mathsf{Eq}(a(v_i,\nu(m)),\nu(n))\). Tag 14 yields
\(\mathsf{Le}_i(\nu(m),\nu(n))\). For the positive strict comparison,
put \(d=L(n,m)>0\); the arithmetic equations give \(A(m,d)=n\) and
\(S(D(d))=d\). Thus \(R_+(m,d)\) has root
\(\mathsf{Eq}(a(\nu(m),s(\nu(D(d)))),\nu(n))\), which is the stated matrix
instance for witness \(\nu(D(d))\). Tag 14 yields its exact strict
comparison root. A numeral has no variables, so these substitutions do not
alter the supplied binder index. These raw trees satisfy \(\mathsf{Safe}\).

Now verify the negative weak comparison with \(m>n\). Its major assumption
\(H\) has root \(C=\exists x_i(x_i+\bar m=\bar n)\), marker 1.
Its minor assumption \(K\) is exactly
\(x_y+\bar m=\bar n\), marker 2, where \(y=i+3\).
The right-addition recurrence proves
\(x_y+\bar m=S^m x_y\) as an intermediate open formula. Symmetry and
transitivity with \(K\) give \(S^m x_y=\bar n\), exactly the root
required by \(\mathsf{Clash}(v_y,m,n,\cdot)\). That constructor produces
bottom and retains the minor marked assumption once.

For the negative strict comparison \(m\geq n\), the minor is
\(\bar m+Sx_y=\bar n\). The left-addition recurrence proves
\(\bar m+Sx_y=S^m(Sx_y)\). Its last term has exactly the code
\(\mathsf{it}(v_y,Sm)\), by the internal iterate identity already proved.
Symmetry and transitivity therefore give the root needed by
\(\mathsf{Clash}(v_y,Sm,n,\cdot)\), and \(Sm>n\) licenses that
constructor. It again returns bottom with the one marked minor leaf.

For both orientations, the new tag-16 row has parameters precisely
\([i,c,y,1]\), major \(H\), and the minor bottom derivation. Its
formula-specific discharge removes the marker-2 minor instance, nonvacuously.
Every other minor open leaf is an unmarked member of \(\mathcal A_*\).
The eigenvariable \(y=i+3\) differs from \(i\), is at least 3, and is
absent from every such leaf's all-variable set, since those sets are subsets
of \(\{0,1\}\). The major comparison has only the bound index \(i\)
in its all-variable set, and the conclusion bottom has no variables.
The remaining major assumption is therefore also free of every occurrence
of \(y\). These facts verify all three factors and the bounded product
of the selected existential-elimination freshness test, including bound
occurrences. Its open-list order is the filtered minor list followed by the
major list, exactly as in the defined template.

The new tag-9 row, with parameter \([0]\), then removes the major marker-1
assumption \(C\), which remains once. Its conclusion is \(\mathsf{Neg}(C)\) and
its remaining leaves are unmarked \(\mathcal A_*\) leaves. This proves
\(\mathsf{Safe}\) and the sentence-root test for the two negative
comparison records. Previously discharged label-0 assumptions inside copied
\(D_Z,D_L\) records cause no collision: they are absent from those roots'
open lists, and each earlier discharge still sees only its copied branch.

The sentence assertions here follow from the internal variable-table
equations. A numeral has every variable bit zero. In either comparison
matrix the only possible free index is its supplied \(i\); at the
existential constructor the index-\(i\) free bit is multiplied by zero,
and at every other index the matrix bit was already zero. Negation preserves
these zero bits. The type-table cases make every such root a formula, and
the free-index filter is therefore empty. This proves its actual
\(\mathsf{Sent}\) test for all numerical binder indices, including zero.

Every arithmetic fact about \(L,E,\mathsf{lt}\) in this argument is an
internal equation or comparison consequence proved for those descriptions
in \cref{annex:p5.checker-description}, using the PA arithmetic laws from
\cref{annex:p5.internal-pr-graphs}. The conditions selecting the displayed
branches are therefore proved to make their success tests one. The zero
failure branches remain total without any claimed proof on their failing
domain.

Finally, \(R_{\neg\bot}\) has exactly two rows. Its first row is the
marked bottom assumption, and its second row is tag 9 with that row as its
sole premise, formula \(\mathsf{Neg}(\mathsf{bc})\), parameter list \([0]\), and
zero marker. Its discharge removes the unique bottom leaf. The first row
has one parent, the last zero, and its final open list is empty. Thus it
passes the selected checker after packing. At equal comparison endpoints,
the positive weak comparison uses witness zero, while the negative strict
comparison leaves a positive successor after peeling. In particular these
arguments cover \(m=n=0\); neither an absent predecessor of zero nor a
nonempty peeling range is assumed. All the tabled proof assertions now
follow by packing the verified sentence-root trees.
\end{proof}
\end{annexitem}

\begin{annexitem}{annex:p5.fixed-term-literal-builders}{thm:p05-fixed-term-literal-builders}{Fixed formation, arithmetic value and exact literal transport}
\begin{proof}
First, the descriptions in the body are finite in the required sense.
Each row template is a fixed composition of earlier descriptions. The four
arithmetic recurrences and the peeling recurrence use the declared last
argument; their step descriptions call only earlier descriptions and their
supplied preceding state. The peeling bound \(b\) is supplied before its
run. Its branch tests, and the sign/ordering tests selecting literal builders,
are the previously fixed Boolean descriptions. The two derived induction
records are fixed finite compositions. An instance builder performs exactly
one or two declared universal eliminations. Consequently there is no cyclic
description dependency or proof-search call. Even a failure branch has its
specified numerical value. The graph theorem supplies unique outputs for
each of these fixed descriptions.

For a fixed term, metainduction on its finite formation assigns the
descriptions \(t^\#,t^\nu,R_t\). The induction assertion is an actual
\(\PA\) theorem, with all its numerical tuple variables universally
quantified, asserting \(\mathsf{Safe}\) with root
\(\mathsf{Eq}(t^\nu,\nu(t^\#))\). The zero and variable cases are exactly
their tag-18 identity rows. In the successor case, congruence lifts the
child equality and \(\nu(Sk)=s(\nu(k))\) identifies its new right term.
For addition, the first congruence replaces \(u^\nu\) in the left
argument while keeping \(v^\nu\). The second replaces \(v^\nu\) in
the right argument while keeping the already normalized left numeral.
Their middle term is literally the same code. Transitivity gives
\[
 \mathsf{Eq}(a(u^\nu,v^\nu),a(\nu(u^\#),\nu(v^\#))).
\]
The previously verified addition proof supplies its last normalization
step. Two transitivity joins yield the required output. Multiplication
uses the same two child equalities in the specified multiplication contexts,
then the verified multiplication proof. These steps retain only the fixed
unmarked axiom leaves. This is a finite metainduction producing internal
PA proofs for the separately fixed descriptions; it does not quantify over
term descriptions inside \(\PA\).

The same formation induction using the syntax equations proves that
\(t^\nu\) is a variable-free term and is the literal numeral substitution
instance of the original term. A numeral insertion is returned without
being traversed as a new input, and the no-conflict table commutes with all
of its term constructors. Thus the assertion covers any order of the
declared distinct free-variable list; the numerical tuple positions were
fixed before the construction.

For the arithmetic antecedents, \(\PA\) proves that the fixed graph
\(A\) agrees with object-language addition, as already established in
\cref{ex:p05-pa-pair-cons-graphs}. Induction on the second argument of
\(M\) similarly gives agreement with object-language multiplication:
the graph's zero equation and \(q_6\) both give zero, and its successor
equation, the addition identification and \(q_7\) give the same next
value. Uniqueness proves both directions. The initial zero/successor
descriptions already have their literal equations. Metainduction on the
fixed term therefore gives, inside \(\PA\), equality of its value as
an object-language term at \(\bar m\) and the graph value \(t^\#(\bar m)\).
This is an arithmetic equality, not a uniform semantic evaluation theorem.
Equality reasoning, and the fixed inherited comparison definitions with the
internal PA order laws, consequently turn a true fixed signed literal
antecedent \(\lambda(\bar m)\) into exactly the numerical condition in
the normalized-literal table. For negated comparisons the same equivalences
and trichotomy give the indicated complementary condition.

Assume that condition and take the corresponding normalized raw proof
\(Q\). Its root is \(W(\nu(u^\#),\nu(v^\#))\). The chosen hole
index \(j\) is larger than the comparison binder \(i\), if present,
and absent from every term argument, since those are variable-free. It
is therefore a permissible hole in both displayed formula contexts. The
first symmetry proof has equality root \(\mathsf{Eq}(\nu(u^\#),u^\nu)\).
The two substitution equations for its equality-elimination matrix are
\[
\begin{aligned}
 W(v_j,\nu(v^\#))[\nu(u^\#)/x_j]
       &=W(\nu(u^\#),\nu(v^\#)),\\
 W(v_j,\nu(v^\#))[u^\nu/x_j]
       &=W(u^\nu,\nu(v^\#)).
\end{aligned}
\tag{9}
\]
Brackets in (9) mean the actual selected substitution graph. The new
no-conflict theorem proves these identities also when \(W\) has an
existential binder and an outer negation: the replacement terms have no
variables, and that binder differs from the hole. The first tag-19 join
therefore has precisely the latter root. The second join, using symmetry
of the \(v\) evaluation proof and matrix \(W(u^\nu,v_j)\), gives
\(W(u^\nu,v^\nu)\) by the corresponding two identities. It adds no
open assumptions except the already allowed unmarked leaves of its three
input trees. Both joins have the actual parameter list order
\([\text{left term},\text{right term},\text{matrix},j]\).

The final root is a sentence. It is also the literal numeral-instance code
\(\lambda^\nu(\bar m)\). For a comparison, the existential binder was
the least index absent from the original two terms when that abbreviation
was expanded. Those indices and that code were fixed at that stage.
Numeral substitution does not rename this binder, and the construction
retains its supplied index \(i\), even if index zero becomes available
after substitution. Thus there is no silent replacement of the target by
an alpha-equivalent code. For a signed equality there is no comparison
binder to retain. The outer negation, when present, is likewise preserved
literally by (9).

The resulting \(\mathsf{Safe}\) assertion and its sentence root give
the selected \(\mathsf{Prf}^{19}_{\PA}\) formula after packing. The
success branch is the specified output of \(d_\lambda\). Expand all
its finitely many graph applications, existentially introduce the unique
output \(p\), and apply implication introduction and universal closure
in the tuple parameters. This is exactly the internal formula displayed
in the body. The negative-bottom case uses its explicit two-row proof;
for positive bottom, bottom elimination gives the displayed implication
with contradictory antecedent, and no successful output of its zero-valued
builder is asserted.

An empty parameter list causes no exceptional proof step: the description
is nullary, the tuple quantifier prefix is empty, and the same fixed
formation argument applies. Numeral zero, multiplication by zero, witness
zero for a weak comparison, a zero peeling bound, equal endpoints and empty
identity-proof contexts were included in the preceding base cases. This
completes the supplied literal constructors and their internal invariants.
The bounded connective and quantifier constructors are the next dependency;
no evidence-to-proof or positive-introspection claim follows from literals
alone.
\end{proof}
\end{annexitem}
